%!TEX root = ../article.tex

The field $$K = \Rational(\lambda_0) = \Rational\left(i\sqrt{\frac{-1+\sqrt{13}}{2}}\right)$$ is a degree $4$ non-Galois extension of $\Rational$.  Its Galois closure $L$ has degree $8$ and has Galois group the dihedral group with $8$ elements which we think of as the automorphisms of a fixed square.  The field $K$ is the fixed points of the subgroup generated by a reflection in one of the sides in the square.  In particular, the orbit of $K$ under the Galois group $\text{Gal}(L/\Rational)$ consists of two fields, $K$ and $$K' = \Rational\left(i\sqrt{\frac{-1-\sqrt{13}}{2}}\right).$$  Each of $K$ and $K'$ is fixed pointwise by two elements of $\text{Gal}(L/\Rational)$ and each has a single nontrivial automorphism over $\Rational$.  The field automorphism of $K/\Rational$ acts on $\cH_0^p$ by the diagram automorphism which interchanges $P$ and $Q$.  In particular, using the Galois action we can only construct a single new fusion category which is defined over $K'$.  This category is non-unitary and the dimensions of the objects are given by replacing $\sqrt{13}$ with $-\sqrt{13}$ everywhere.

The story is similar for $ \Rational(\lambda_1)$ and $\cH_1^p$. Again, there is a single nontrivial automorphism of the field which acts on the fusion category by diagram automorphism.  There are two Galois conjugate fusion categories which are non-unitary and whose dimensions are given by the action of $\text{Gal}(\Rational(d_1^2)/\Rational)$ on the old dimensions.